% ======================================================================
%  BIBM / IEEE 课程论文模板（支持中英文混排）
%  适用课程：《生物文本挖掘与知识发现》等
%  编译方式：pdflatex  ->  bibtex(可选)  ->  pdflatex  ->  pdflatex
%  依赖宏包：IEEEtran.cls（BIBM 官方模板，请与本文件放同一目录）
%
%  本文档内容：论文标题 + 主要研究设想
%  详细论文方案见同目录下《论文方案.md》
% ======================================================================
\documentclass[conference]{IEEEtran}
\IEEEoverridecommandlockouts

% ---------- 基础宏包 ----------
\usepackage{cite}
\usepackage{amsmath,amssymb,amsfonts}
\usepackage{algorithmic}
\usepackage{graphicx}
\usepackage{textcomp}
\usepackage{xcolor}
\usepackage{url}
\usepackage{booktabs}   % 更漂亮的三线表（可选）

% ---------- 页码设置 ----------
% IEEE 模板默认不显示页码；通过 fancyhdr 在页脚居中显示页码。
\usepackage{fancyhdr}
\pagestyle{fancy}
\fancyhf{}                    % 清空默认页眉页脚
\fancyfoot[C]{\thepage}       % 页脚居中显示页码
\renewcommand{\headrulewidth}{0pt}  % 去掉页眉横线
\renewcommand{\footrulewidth}{0pt}  % 去掉页脚横线

% ---------- 页脚左下角灰色水印 ----------
% 课程论文声明：Agent 生成内容提示，灰色小字，左下角紧凑三行。
% 若不需要，注释掉下面 \fancyfoot[L]{...} 整段即可。
\fancyfoot[L]{%
  \scriptsize\color{gray}%
  \begin{tabular}[b]{@{}l@{}}
    课程论文支持使用 Agent 生成内容，\\
    这只是起点。人工内容审查\\
    和驾驭才是终点。%
  \end{tabular}%
}
\setlength{\footskip}{24pt}   % 给页脚留出足够空间，避免与正文重叠

% IEEEtran 的首页 (\maketitle) 默认使用 plain 样式，会覆盖上面的 fancy 设置。
% 这里把 plain 和 IEEE 标题页样式也改成同样的页脚，保证首页同样显示页码和水印。
\fancypagestyle{plain}{%
  \fancyhf{}%
  \fancyfoot[C]{\thepage}%
  \fancyfoot[L]{%
    \scriptsize\color{gray}%
    \begin{tabular}[b]{@{}l@{}}
      课程论文支持使用 Agent 生成内容，\\
      这只是开始。人工内容审查\\
      和驾驭才是终点。%
    \end{tabular}%
  }%
  \renewcommand{\headrulewidth}{0pt}%
  \renewcommand{\footrulewidth}{0pt}%
}
\fancypagestyle{IEEEtitlepagestyle}{%
  \fancyhf{}%
  \fancyfoot[C]{\thepage}%
  \fancyfoot[L]{%
    \scriptsize\color{gray}%
    \begin{tabular}[b]{@{}l@{}}
      课程论文支持使用 Agent 生成内容，\\
      这只是开始。人工内容审查\\
      和驾驭才是终点。%
    \end{tabular}%
  }%
  \renewcommand{\headrulewidth}{0pt}%
  \renewcommand{\footrulewidth}{0pt}%
}

% ---------- 行号设置 ----------
% 在页面左侧显示行号，方便老师批改/审稿时定位。
% 若提交最终版不需要行号，注释掉 \pagewiselinenumbers 一行即可。
\usepackage[switch]{lineno}   % switch: 双栏时行号显示在外侧
\pagewiselinenumbers          % 每页行号从 1 重新开始（也可改用 \linenumbers 全文连续编号）

% ---------- 中文支持（pdflatex 下也能用） ----------
% 说明：
%   1) 使用 ctex 宏包并开启 UTF8，即可在 pdflatex 下输入中文；
%   2) 若你想用 xelatex，可去掉下面这行，改用 \usepackage[UTF8]{ctex} 并设置中文字体。
\usepackage[UTF8]{ctex}

\def\BibTeX{{\rm B\kern-.05em{\sc i\kern-.025em b}\kern-.08em
    T\kern-.1667em\lower.7ex\hbox{E}\kern-.125emX}}

\begin{document}

% ======================================================================
%  标题 / 作者 / 单位  ——  按需修改
% ======================================================================
\title{From Co-occurrence to Signal: Auditing and Correcting
Research-Attention Bias in Cross-Kingdom Literature-Mined
Gene--Trait Knowledge Bases\\
\large 从共现到信号：跨物种文献挖掘基因--性状知识库中
研究注意力偏差的审计与校正}

\author{
\IEEEauthorblockN{陈佳伸}
\IEEEauthorblockA{
\textit{2023317220404} \\
\textit{人工智能2302}
邮箱: smiletrashbag@webmail.hzau.edu.cn}
}

\maketitle

% ======================================================================
%  摘要 / 关键词
% ======================================================================
\begin{abstract}
文献挖掘型知识库正被当作事实性资源广泛复用，但其关联强度的底层是共现计数，
而共现同时混入了\textbf{生物学效应}与\textbf{研究注意力}两个来源。
本文提出一项跨物种研究设想：以 RiceMind（水稻）与 CancerAlterome（人类癌症）
这两个出自同一实验室、共享同一「共现抽取 + 本体归一化」范式的资源为对象，
用库外文献基线为基因--性状对建立期望共现模型，
并以残差作为注意力校正后的信号。
设想的关键在于一个方法学处理：经查证，RiceMind 的 Tier~2/Tier~3 分层
本身即由文章数阈值定义（$>10$ 与 $\leq 10$ 篇），
故基于该分层验证任何基于计数的校正都将构成循环；
本文据此仅以与文献量无关的 Tier~1 作为标签，并对比较按文章数分层。
预期贡献为：一套可移植的校正协议、对分层循环性的显式处理，
以及首次对文献挖掘知识库注意力偏差的跨物种量化。
\end{abstract}

\begin{IEEEkeywords}
文献挖掘, 知识库偏差, 共现归一化, 发表偏倚,
本体归一化, 跨物种迁移
\end{IEEEkeywords}

% ======================================================================
%  I. 研究背景与问题
% ======================================================================
\section{研究背景与问题 Background}

生物医学文献的规模已远超人工阅读能力，文献挖掘因此成为构建结构化知识的主要手段。
Cancer-Alterome 与 RiceMind
是该范式下两个成熟的资源：前者面向人类癌症，
将\textbf{遗传改变}$\rightarrow$\textbf{下游调控事件}归一化到 GO/HPO；
后者面向水稻，将基因--性状关联归一化到 GO/TO/PO/RTO。
两者出自同一实验室（HZAU-BioNLP Lab），共享同一套
「共现抽取 + 本体归一化 + 来源分级」的方法论，区别仅在于物种与表型空间。

这类资源被下游直接当作「已知事实」使用——用于候选基因排序、
通路富集、以及作为机器学习模型的弱标签。然而其关联强度的底层量是
\textbf{文献共现计数}。一个基因--性状对共现得多，可能因为该关联在生物学上强，
也可能仅仅因为\textbf{这个基因和这个性状都被研究得多}。
后者即\textbf{研究注意力}，它与生物学效应在计数中被不可分辨地混合。

这一偏差会带来一个具有实际后果的反馈回路：
若知识库被不加校正地用于下游，其输出将系统性偏向已被充分研究的实体，
从而\textbf{放大而非缩小}现有研究格局的不均衡——
研究得多的基因被预测为更重要，进而获得更多研究。
在作物育种与癌症靶点发现两个场景中，这意味着
前者会忽视地方品种中未被充分研究的耐逆基因，
后者会重复聚焦于已饱和的靶点。

现有资源普遍以置信度分层应对该问题，但这些分层并未与文献量解耦。
经查证 RiceMind 的官方文档，其 Tier~2 与 Tier~3 的划分标准分别是
「$>10$ 篇独立文章」与「$\leq 10$ 篇文章」%
\footnote{\url{http://lit-evi.hzau.edu.cn/ricemind-webpage/dataset/}}，
\textbf{即某一层的定义本身就是文章数的函数}。
这一事实尚未被既有文献指出，却直接影响下游对分层的使用方式。

\textbf{本文要回答的三个问题：}
\begin{itemize}
  \item \textbf{RQ1（量化）} 研究注意力对关联强度的影响有多大？
        它与生物学信号在多大程度上可分离？
  \item \textbf{RQ2（校正）} 用外部文献基线对共现做归一化，
        能否得到比原始计数更接近生物学真值的排序？
  \item \textbf{RQ3（迁移）} 在一个物种上确定的校正，
        能否无损迁移到另一物种及其表型空间？
\end{itemize}

% ======================================================================
%  II. 主要研究设想
% ======================================================================
\section{主要研究设想 Main Research Idea}

\textbf{核心设想：把「超出研究注意力所能解释的部分」定义为关联信号。}

设基因 $g$ 与性状（或癌症过程）$t$ 的库内共现计数为 $n(g,t)$，
库内边缘频次为 $n_g$、$n_t$，总量为 $N$；
另设 $m_g$、$m_t$ 为 PubMed 中 $g$ 与 $t$ 各自的独立文献量，$M$ 为检索空间总量。
则\textbf{仅由注意力驱动}的期望共现为
\begin{equation}
E_{\text{ext}}[n(g,t)] = N \cdot \frac{m_g}{M} \cdot \frac{m_t}{M}
\end{equation}
残差
\begin{equation}
r(g,t) = n(g,t) - E_{\text{ext}}[n(g,t)]
\end{equation}
即为本文所定义的注意力校正后信号。

设想的要点不在于提出新的归一化公式，
而在于主张\textbf{期望值的边缘频次应取自库外}：
库内计数本身即是被污染的观测量，
以它构造期望无法区分「该基因被研究得多」与「该关联真实且强」。

\textbf{由此形成一个可直接实施的实验设计：}
在同一测试集上比较三种排序——
原始计数 $n$、仅用库内边缘频次的 $\mathrm{PMI}_{\text{int}}$、
以及用库外基线的 $\mathrm{PMI}_{\text{ext}}$。
其中 $\mathrm{PMI}_{\text{int}}$ 与 $\mathrm{PMI}_{\text{ext}}$ 的
\textbf{排名差异即构成文献注意力偏差的直接测量值}，
该差值在既有文献中尚未被报告。

\textbf{验证方案的立意：规避分层的循环性。}
如第一节所述，RiceMind 的 Tier~2/Tier~3 由文章数阈值定义，
故基于该二分层的验证是循环的，这是本设想在方法学上必须处理的核心障碍。
应对方式有三：
（i）\textbf{标签选择}上仅使用 \textbf{Tier~1} 作为正例——
其定义依据是人工审编来源（Oryzabase/RAP-DB）或实验证据码（EXP/IDA/IMP），
\textbf{与文献量无关}，故可作为干净标签，而 Tier~2/3 仅作分层变量；
（ii）\textbf{按文章数区间分层}（如 $1$--$5$、$6$--$10$、$11$--$50$、$>50$），
在各区间\textbf{内部}比较排序质量，
以控制「Tier~1 基因本来就被研究得多」这一残余混淆；
（iii）以\textbf{跨物种迁移}作为独立检验——
在水稻侧确定的校正参数直接应用于癌症侧，
该检验不依赖任何一方的分层体系，因而不受循环性影响。

\textbf{待检验假设：}
\begin{itemize}
  \item \textbf{H1（校正有效性）} 分层条件下，
        $\mathrm{PMI}_{\text{ext}}$ 对 Tier~1 关联的 $P@k$
        显著高于原始计数且不低于 $\mathrm{PMI}_{\text{int}}$；
  \item \textbf{H2（偏差可量化）} 两种 PMI 的排名差异与实体级文献量显著相关，
        且该相关在两库中具有可比的函数形式；
  \item \textbf{H3（跨物种迁移）} 水稻侧校正参数应用于癌症侧时排序质量无显著下降；
  \item \textbf{H4（漏检回收）} 外部基线驱动的残差排序能提升
        「高文献量、零库内记录」实体的关联召回率。
\end{itemize}

% ======================================================================
%  III. 实验设计
% ======================================================================
\section{实验设计 Experimental Design}

本节给出检验 H1--H4 的具体方案，涵盖数据构建、标签设定、
对比方法、实验矩阵、评价指标与稳健性检验。

\subsection{数据构建 Data Construction}

三个数据面分别对应式 (1)(2) 所需的量：

\begin{itemize}
  \item \textbf{库内关联与计数}（$n(g,t)$、$n_g$、$n_t$、$N$）：
        RiceMind 经 MCP 服务取全部 GTA 记录，
        字段含 \texttt{ontology\_id}、\texttt{confidence\_tier}、
        \texttt{evidence\_codes}、\texttt{article\_count}；
        CancerAlterome 经 \texttt{gene\_var\_table\_update\_paginator/}
        与 \texttt{\_by\_cancer/} 两个端点按基因、按癌症分别拉取。
        \textbf{须避开}无 \texttt{\_paginator} 后缀的端点——
        实测其静默忽略查询参数并返回固定的 315 条 ERBB2 数据，
        返回码仍为 HTTP 200。
  \item \textbf{库外文献基线}（$m_g$、$m_t$、$M$）：
        经 NCBI E-utilities \texttt{esearch.fcgi} 对每个实体检索 PubMed，
        记录命中数。检索式须预先固定并随代码公开。
  \item \textbf{实体对齐}：RiceMind 侧统一到 RAP ID，
        CancerAlterome 侧沿用其基因符号。
        两库之间\textbf{不做}基因层面的同源映射；
        跨物种迁移仅在\textbf{分布层面}（相关参数与函数形式）进行。
\end{itemize}

\textbf{评估单元}为三元组 $(g,t,\text{库})$。
两库实测规模：RiceMind 基因 18{,}983 个、性状 5{,}371 条；
CancerAlterome 基因 5{,}474 个。

\subsection{标签设定：正--未标注框架 Label Construction}

\textbf{本研究没有可靠的负例，这是实验设计的首要约束。}

「某基因与某性状不存在关联」无法从库中确证：
库中未出现既可能是关联不存在，也可能是文献未被挖掘。
若把「非 Tier~1」直接当作负例，
则假阴性会被计入假阳性，从而系统性低估校正方法的收益。

据此采用\textbf{正--未标注}（positive--unlabeled, PU）设定：

\begin{itemize}
  \item \textbf{正例}：Tier~1 关联。其判据为人工审编来源
        （Oryzabase/RAP-DB）或实验证据码（EXP/IDA/IMP），
        \textbf{与文献量无关}。
  \item \textbf{未标注}：其余全部关联。
        Tier~2/3 \textbf{不参与}正例或负例的构造，
        仅作为\textbf{分层变量}使用。
\end{itemize}

该设定直接规避了第一节指出的循环性：
由于 Tier~2/3 由文章数阈值定义，
任何以「Tier~1 对 Tier~3」为二分类标签的评测，
其难度都部分等价于「模型能否数到 10」。

\subsection{对比方法 Baselines}

\begin{table}[h]
  \centering
  \small
  \caption{对比方法。B4 为循环性对照，用于实证分层的循环性，不是竞争方法。}
  \label{tab:baselines}
  \begin{tabular}{lll}
    \toprule
    编号 & 方法 & 说明 \\
    \midrule
    B1 & 原始计数 $n$ & 现状基线 \\
    B2 & $\mathrm{PMI}_{\text{int}}$ & 仅校正库内密度 \\
    B3 & $\mathrm{PMI}_{\text{ext}}$ & 本文主张 \\
    \midrule
    B4 & \texttt{article\_count} & 循环性对照 \\
    \bottomrule
  \end{tabular}
\end{table}

\textbf{B4 的作用。}由于 RiceMind 的 Tier~2/3 边界本身即文章数阈值，
仅用 \texttt{article\_count} 单一变量预测「Tier~2 对 Tier~3」
即可取得接近完美的表现。B4 不参与方法优劣的比较，
而是\textbf{作为循环性的实证演示}：
若 B4 在该任务上显著优于全部真实方法，
则说明该任务测的是文献量而非知识质量；
据此可进一步量化「既有分层作为评测基准的失效程度」。

\subsection{实验矩阵 Experiment Matrix}

\begin{table}[h]
  \centering
  \small
  \caption{实验矩阵。E3 为方法学对照，其余四项分别对应 H1--H4。}
  \label{tab:matrix}
  \begin{tabular}{llll}
    \toprule
    编号 & 目的 & 对比 & 指标 \\
    \midrule
    E1 & 偏差量化 & B1/B2/B3 & $\Delta$rank 对 $\log m_g$ \\
    E2 & 校正有效性 & B1/B2/B3 & 分层 $P@k$、$EF@k$ \\
    E3 & 循环性实证 & B4 对其余 & Tier~2/3 上 $P@k$ \\
    E4 & 跨物种迁移 & B3 跨库应用 & 非劣性检验 \\
    E5 & 漏检回收 & B1 对 B3 & Recall@$k$ \\
    \bottomrule
  \end{tabular}
\end{table}

\textbf{E1（对应 H2）}在全部评估单元上计算
$\mathrm{PMI}_{\text{int}}$ 与 $\mathrm{PMI}_{\text{ext}}$ 的排名差
$\Delta\text{rank}$，并对实体级文献量 $\log m_g$ 做回归，
比较两库所得函数的形状与量级。
\textbf{E2（对应 H1）}在各文章数区间\textbf{内部}比较 B1--B3 对 Tier~1 的排序质量。
\textbf{E3}单独运行，其结论不用于评判方法，而用于说明分层作为基准的局限性。
\textbf{E4（对应 H3）}把水稻侧确定的参数直接施加于癌症侧，做非劣性检验。
\textbf{E5（对应 H4）}构造「高 PubMed 文献量、库内记录为零或极少」的实体集，
比较残差排序对相关 Tier~1 关联的召回。

\subsection{评价指标 Metrics}

由于负例不可靠，主指标选用不依赖真负例的排序指标：

\begin{itemize}
  \item $P@k$：前 $k$ 位中 Tier~1 的比例，$k\in\{10,25,50,100\}$；
  \item 富集倍数 $EF@k$：前 $k$ 位 Tier~1 比例除以全库 Tier~1 比例；
  \item Recall@$k$：前 $k$ 位覆盖的 Tier~1 占全部 Tier~1 的比例；
  \item Kendall $\tau$：校正前后排名的整体漂移幅度，
        用于 E1 的偏差度量，\textbf{不}用作准确率。
\end{itemize}

\textbf{不以 AUC-ROC 作主指标。}在 PU 设定下真负例集不可知，
AUC 的数值依赖人为构造负例的方式，跨实验不可比。

\subsection{统计检验与稳健性 Robustness}

\begin{itemize}
  \item \textbf{置信区间}：对评估单元做 $B=1000$ 次自助重采样，
        报告 95\% 百分位置信区间。
  \item \textbf{配对检验}：在同一评估单元上比较 B1 与 B3，
        使用配对自助法差值检验，不依赖正态假设。
  \item \textbf{非劣性检验（E4）}：预设边界 $\delta$，
        若跨物种应用后 $P@k$ 下降不超过 $\delta$ 则判定迁移成功。
        $\delta$ 须在实验前固定并说明依据，\textbf{不得事后调整}。
  \item \textbf{敏感性分析}：分别改变
        （i）PubMed 检索式的同义词扩展策略，
        （ii）分层区间边界，
        （iii）$k$ 的取值，
        观察结论是否稳定。
  \item \textbf{多重比较}：跨库、跨 $k$、跨分层区间构成多重比较，
        报告 Benjamini--Hochberg 校正后的 $q$ 值。
\end{itemize}

\subsection{可复现性 Reproducibility}

公开两库的接口版本与访问日期、全部 PubMed 检索式、
分层边界、随机种子，以及按端点原始返回解析的中间数据。
所有计数均可由公开接口重新取得，不依赖私有快照。

\textbf{说明：}本文为研究设想，上述 H1--H4 为待检验假设；
本节给出检验方案，本文不预设任何实验结果。

\end{document}
